\documentclass[10pt,a4paper]{amsart}
\usepackage{amsmath,amssymb,amsthm,graphicx}
% Retained minimal preamble for the Stacks Project chapter varieties.tex.
% Verbatim excerpt of the chapter preamble: Xy-pic setup (native preamble lines 14-18), the native theorem
% declarations with the mathematical macros (native preamble lines 154-217), and the project own no-op
% definitions of the reference/slogan/history display environments (native wrapper.tex lines 1-3) that the
% counted statement uses for its slogan. Copied mechanically; nothing added, removed or redefined.
% For commutative diagrams we use Xy-pic
\usepackage[all]{xy}

% We use 2cell for 2-commutative diagrams.
\xyoption{2cell}
% Theorem environments.
%
\theoremstyle{plain}
\newtheorem{theorem}[subsection]{Theorem}
\newtheorem{proposition}[subsection]{Proposition}
\newtheorem{lemma}[subsection]{Lemma}

\theoremstyle{definition}
\newtheorem{definition}[subsection]{Definition}
\newtheorem{example}[subsection]{Example}
\newtheorem{exercise}[subsection]{Exercise}
\newtheorem{situation}[subsection]{Situation}

\theoremstyle{remark}
\newtheorem{remark}[subsection]{Remark}
\newtheorem{remarks}[subsection]{Remarks}

\numberwithin{equation}{subsection}

% Macros
%
\def\lim{\mathop{\mathrm{lim}}\nolimits}
\def\colim{\mathop{\mathrm{colim}}\nolimits}
\def\Spec{\mathop{\mathrm{Spec}}}
\def\Hom{\mathop{\mathrm{Hom}}\nolimits}
\def\Ext{\mathop{\mathrm{Ext}}\nolimits}
\def\SheafHom{\mathop{\mathcal{H}\!\mathit{om}}\nolimits}
\def\SheafExt{\mathop{\mathcal{E}\!\mathit{xt}}\nolimits}
\def\Sch{\mathit{Sch}}
\def\Mor{\mathop{\mathrm{Mor}}\nolimits}
\def\Ob{\mathop{\mathrm{Ob}}\nolimits}
\def\Sh{\mathop{\mathit{Sh}}\nolimits}
\def\NL{\mathop{N\!L}\nolimits}
\def\CH{\mathop{\mathrm{CH}}\nolimits}
\def\proetale{{pro\text{-}\acute{e}tale}}
\def\etale{{\acute{e}tale}}
\def\QCoh{\mathit{QCoh}}
\def\Ker{\mathop{\mathrm{Ker}}}
\def\Im{\mathop{\mathrm{Im}}}
\def\Coker{\mathop{\mathrm{Coker}}}
\def\Coim{\mathop{\mathrm{Coim}}}

% Boxtimes
%
\DeclareMathSymbol{\boxtimes}{\mathbin}{AMSa}{"02}

%
% Macros for moduli stacks/spaces
%
\def\QCohstack{\mathcal{QC}\!\mathit{oh}}
\def\Cohstack{\mathcal{C}\!\mathit{oh}}
\def\Spacesstack{\mathcal{S}\!\mathit{paces}}
\def\Quotfunctor{\mathrm{Quot}}
\def\Hilbfunctor{\mathrm{Hilb}}
\def\Curvesstack{\mathcal{C}\!\mathit{urves}}
\def\Polarizedstack{\mathcal{P}\!\mathit{olarized}}
\def\Complexesstack{\mathcal{C}\!\mathit{omplexes}}
% \Pic is the operator that assigns to X its picard group, usage \Pic(X)
% \Picardstack_{X/B} denotes the Picard stack of X over B
% \Picardfunctor_{X/B} denotes the Picard functor of X over B
\def\Pic{\mathop{\mathrm{Pic}}\nolimits}
\def\Picardstack{\mathcal{P}\!\mathit{ic}}
\def\Picardfunctor{\mathrm{Pic}}
\def\Deformationcategory{\mathcal{D}\!\mathit{ef}}
\newenvironment{reference}{}{}
\newenvironment{slogan}{}{}
\newenvironment{history}{}{}
\usepackage{cleveref}
\crefname{originallemma}{original lemma}{original lemmas}
\Crefname{originallemma}{Original lemma}{Original lemmas}
\makeatletter
\newlabel{lemma-normalize-noetherian-dim-1}{{0C45}{source}{}{sourcepointer-1}{}}
\newlabel{lemma-normalize-noetherian-dim-1@cref}{{[originallemma][0][]0C45}{[1][0][]source}{}{}{}}
\makeatother
\crefname{originalreference}{original reference}{original references}
\Crefname{originalreference}{Original reference}{Original references}
\makeatletter
\newlabel{morphisms-lemma-nagata-normalization}{{035S}{source}{}{sourcepointer-2}{}}
\newlabel{morphisms-lemma-nagata-normalization@cref}{{[originalreference][0][]035S}{[1][0][]source}{}{}{}}
\makeatother
\crefname{originalreference}{original reference}{original references}
\Crefname{originalreference}{Original reference}{Original references}
\makeatletter
\newlabel{more-algebra-proposition-ubiquity-J-2}{{07PJ}{source}{}{sourcepointer-3}{}}
\newlabel{more-algebra-proposition-ubiquity-J-2@cref}{{[originalreference][0][]07PJ}{[1][0][]source}{}{}{}}
\makeatother
\crefname{originallemma}{original lemma}{original lemmas}
\Crefname{originallemma}{Original lemma}{Original lemmas}
\makeatletter
\newlabel{lemma-delta-invariant-is-zero}{{0C3U}{source}{}{sourcepointer-4}{}}
\newlabel{lemma-delta-invariant-is-zero@cref}{{[originallemma][0][]0C3U}{[1][0][]source}{}{}{}}
\makeatother
\crefname{originallemma}{original lemma}{original lemmas}
\Crefname{originallemma}{Original lemma}{Original lemmas}
\makeatletter
\newlabel{lemma-normalization-same-after-completion}{{0C3V}{source}{}{sourcepointer-5}{}}
\newlabel{lemma-normalization-same-after-completion@cref}{{[originallemma][0][]0C3V}{[1][0][]source}{}{}{}}
\makeatother
\title{Nine structural facts around the normalization of a reduced Nagata scheme of dimension 1}
\author{The Stacks Project Authors}
\date{stacks-project:varieties.tex:local-snapshot-20261011; GNU FDL 1.2 or later}
\begin{document}
\maketitle
\noindent\textit{Editorial note.} All delivered mathematical text is the author's own, copied verbatim from the pinned Stacks Project chapter file varieties.tex; no statement or proof step was rewritten, reordered or removed.\par\medskip
\noindent\textit{Editorial note.} The proof explicitly reuses the delivered normalization lemma for the stalk statements and the delivered completion lemma for assertions (6)-(9); the short exact sequence and the length identification are carried out in the span.\par\medskip
\noindent\textit{Allowed standard prerequisites (editorial).}\begin{itemize}
\item Normalization of a Noetherian dimension-1 scheme. delivered companion lemma-normalize-noetherian-dim-1. Source of the stalk descriptions and the product-of-Dedekind structure.
\item Finiteness of normalization for Nagata schemes. Morphisms, Lemma morphisms-lemma-nagata-normalization. Finiteness of nu.
\item The regular locus of a reduced Nagata scheme is dense. More on Algebra, Proposition more-algebra-proposition-ubiquity-J-2. Density of the isomorphism locus of nu.
\item Delta invariant vanishes exactly at regular points. delivered companion lemma-delta-invariant-is-zero. Singular-point dichotomy for Q\_x.
\item Integral closures after henselization and completion. delivered companion lemma-normalization-same-after-completion. Assertions (6)-(9) on A-hat and (A')-hat.
\end{itemize}\medskip
\ifdefined\hypertarget\hypertarget{sourcepointer-1}{}\fi
\noindent\textit{Source reference.} Original lemma 0C45: sources/stacks-chow/tags.txt, source lines 15215--15215. Cross-chapter standard prerequisite cited by the author inside the counted span; it is used as a pointer and its own proof is not part of this item.\par\medskip
\ifdefined\hypertarget\hypertarget{sourcepointer-2}{}\fi
\noindent\textit{Source reference.} Original reference 035S: sources/stacks-chow/tags.txt, source lines 4139--4139. Cross-chapter standard prerequisite cited by the author inside the counted span; it is used as a pointer and its own proof is not part of this item.\par\medskip
\ifdefined\hypertarget\hypertarget{sourcepointer-3}{}\fi
\noindent\textit{Source reference.} Original reference 07PJ: sources/stacks-chow/tags.txt, source lines 9799--9799. Cross-chapter standard prerequisite cited by the author inside the counted span; it is used as a pointer and its own proof is not part of this item.\par\medskip
\ifdefined\hypertarget\hypertarget{sourcepointer-4}{}\fi
\noindent\textit{Source reference.} Original lemma 0C3U: sources/stacks-chow/tags.txt, source lines 15204--15204. Cross-chapter standard prerequisite cited by the author inside the counted span; it is used as a pointer and its own proof is not part of this item.\par\medskip
\ifdefined\hypertarget\hypertarget{sourcepointer-5}{}\fi
\noindent\textit{Source reference.} Original lemma 0C3V: sources/stacks-chow/tags.txt, source lines 15205--15205. Cross-chapter standard prerequisite cited by the author inside the counted span; it is used as a pointer and its own proof is not part of this item.\par\medskip
\begin{lemma}
\label{lemma-prepare-delta-invariant}
Let $X$ be a reduced Nagata scheme of dimension $1$. Let $\nu : X^\nu \to X$
be the normalization. Let $x \in X$ denote a closed point. Then
\begin{enumerate}
\item $\nu : X^\nu \to X$ is finite, surjective, and birational,
\item $\mathcal{O}_X \subset \nu_*\mathcal{O}_{X^\nu}$ and
$\nu_*\mathcal{O}_{X^\nu}/\mathcal{O}_X$ is a direct sum of
skyscraper sheaves with value $\mathcal{Q}_x$ at $x$ which is nonzero
if and only if $x$ is a singular point of $X$,
\item $A' = (\nu_*\mathcal{O}_{X^\nu})_x$ is the integral closure
of $A = \mathcal{O}_{X, x}$ in its total ring of fractions,
\item $\mathcal{Q}_x = A'/A$ has finite length equal to the
$\delta$-invariant of $X$ at $x$,
\item $A'$ is a semi-local ring which is a finite product of
Dedekind domains,
\item $A^\wedge$ is a reduced Noetherian complete
local ring of dimension $1$,
\item $(A')^\wedge$ is the integral closure of $A^\wedge$
in its total ring of fractions,
\item $(A')^\wedge$ is a finite product of
complete discrete valuation rings, and
\item $A'/A \cong (A')^\wedge/A^\wedge$.
\end{enumerate}
\end{lemma}

\begin{proof}
We may and will use all the results of
Lemma \ref{lemma-normalize-noetherian-dim-1}.
Finiteness of $\nu$ follows from
Morphisms, Lemma \ref{morphisms-lemma-nagata-normalization}.
Since $X$ is reduced, Nagata, of dimension $1$, we see that the
regular locus is a dense open $U \subset X$ by
More on Algebra, Proposition \ref{more-algebra-proposition-ubiquity-J-2}.
Since a regular scheme is normal, this shows that $\nu$
is an isomorphism over $U$.
Since $\dim(X) \leq 1$ this implies that $\nu$ is not an isomorphism
over a discrete set of closed points $x \in X$.
In particular we see that we have a short exact sequence
$$
0 \to \mathcal{O}_X \to \nu_*\mathcal{O}_{X^\nu} \to
\bigoplus\nolimits_{x \in X \setminus U} \mathcal{Q}_x \to 0
$$
As we have the description of the stalks of $\nu_*\mathcal{O}_{X^\nu}$
by Lemma \ref{lemma-normalize-noetherian-dim-1}, we conclude that
$Q_x = A'/A$ indeed has length equal to the $\delta$-invariant of $X$ at $x$.
Note that $Q_x \not = 0$ exactly when $x$ is a singular point for
example by Lemma \ref{lemma-delta-invariant-is-zero}.
The description of $A'$ as a product of semi-local Dedekind domains follows
from Lemma \ref{lemma-normalize-noetherian-dim-1} as well.
The relationship between $A$, $A'$, and $(A')^\wedge$
we have see in Lemma \ref{lemma-normalization-same-after-completion}
(and its proof).
\end{proof}

\end{document}
